\section{ADR-0001: Product Focus}

\subsection*{Status}

Accepted

\subsection*{Context}

The initial vision considered building a general-purpose engineering memory platform capable of supporting multiple engineering disciplines.

Early product discussions concluded that such a scope would be too broad for an initial implementation and would significantly increase the complexity of product validation.

A narrower domain would enable better understanding of user workflows, more accurate validation, and faster iteration.

\subsection*{Decision}

The first version of Engineering Shadow shall exclusively target RTL Design Engineers and FPGA Design Engineers working on long-term digital hardware design projects.

The product shall be designed around real RTL engineering workflows, with future expansion to other engineering domains considered only after successful validation.

\subsection*{Rationale}

Restricting the initial scope provides several advantages.

\begin{itemize}
    \item The engineering workflow is well understood through first-hand experience.
    \item Product decisions can be validated against real engineering problems.
    \item The system can be optimized for RTL-specific terminology, workflows, and constraints.
    \item Product complexity is significantly reduced during early development.
\end{itemize}

\subsection*{Alternatives Considered}

\begin{enumerate}
    \item Generic engineering assistant.
    \item AI assistant for all software and hardware engineers.
    \item Research assistant independent of engineering domain.
\end{enumerate}

These alternatives were rejected because they dilute the product focus and make early validation significantly more difficult.

\subsection*{Consequences}

All subsequent product decisions shall prioritize the needs of RTL engineers.

Future expansion beyond RTL engineering will be considered only after achieving a successful and validated product within the initial target domain.